% Template for ICASSP-2026 paper; to be used with:
%          spconf.sty  - ICASSP/ICIP LaTeX style file, and
%          IEEEbib.bst - IEEE bibliography style file.
% --------------------------------------------------------------------------
% main_2026temp.tex — PathSelect master-draft text poured into the ICASSP 2026
% paper-kit skeleton, for LENGTH ASSESSMENT against the 4(+1 refs-only) limit.
% Compile needs alongside this file: spconf.sty, IEEEbib.bst (from the paper
% kit), refs.bib, figures/fig1_architecture.png.
% Deviations from the kit's \usepackage line, both layout-legal: booktabs
% (table rules only). \pending{} renders bold brackets, no extra package.
% Final submission will use the ICASSP **2027** kit when it goes live.
% --------------------------------------------------------------------------
\documentclass{article}
\usepackage{spconf,amsmath,amssymb,graphicx,hyperref}
\usepackage{booktabs} % added: kit constrains layout/fonts, not packages

% Example definitions.
% --------------------
\def\x{{\mathbf x}}
\def\L{{\cal L}}
\newcommand{\pending}[1]{\textbf{[PENDING: #1]}}
\newcommand{\Fg}{F_{g}} \newcommand{\Fp}{F_{p}}

% Title.
% ------
\title{REMEMBERING WHERE TO LOOK: CONTINUAL HIERARCHICAL EVIDENCE SELECTION FOR WHOLE-SLIDE IMAGES}
%
% Single address.
% ---------------
\name{Po-Nyo (Aaron) Hung, [co-authors], [advisor]\thanks{[Funding acknowledgement — PI to fill.] ICASSP is single-blind; author list stays on the PDF.}}
\address{[Department, University — PI to fill]}
%
\begin{document}
\ninept % standard choice for ICASSP submissions; comment out for 10pt
%
\maketitle
%
\begin{abstract}
Whole-slide image (WSI) diagnosis operates under an evidence budget: a few patches are selected from thousands, and every downstream decision rests on them. When diagnostic tasks arrive sequentially, existing continual learning protects classifiers or representations, leaving a distinct failure unmeasured---a model can keep its classifier intact yet gradually stop looking at the right tissue. To expose this failure, we freeze the vision--language encoders and the prompt-based diagnostic rule, making a hierarchical selection policy the only trainable component; any cross-method difference in performance is then attributable to where the model looks. Continual learning operates on this policy alone: each task trains a small low-rank residual that is folded exactly into one shared selector, yielding task-free, constant-cost inference. Merging alone provides no systematic protection; forgetting is instead mitigated by replaying stored behavioral snapshots---scores and indices, never images---through dual-level selection distillation and a utility-preserving objective. On four sequential TCGA tasks, this raises final class-incremental accuracy by 34.6 points over sequential fine-tuning in every seed and retains two orders of magnitude more of its originally selected evidence.
\end{abstract}
%
\begin{keywords}
Whole-slide images, continual learning, evidence selection, low-rank adaptation, computational pathology
\end{keywords}
%
\section{Introduction}
\label{sec:intro}
A single WSI contains $10^4$--$10^5$ patches, of which a diagnosis typically depends on a handful. Systems therefore acquire evidence under a budget: choose $B$ patches, then decide. In deployment, tasks arrive sequentially---a hospital adds esophageal cases this quarter, renal the next---and retraining from scratch on all past cohorts is rarely feasible. Continual learning (CL) for WSIs has so far inherited its framing from natural-image classification: preserve the \emph{classifier} or the \emph{representation}. This framing obscures a failure mode unique to budgeted pathology pipelines: a model whose classifier is intact can still stop looking at the right tissue. A question thus arises: when tasks arrive sequentially, does a budgeted diagnostic model forget \emph{where to look}---and can this forgetting be measured, and mitigated, at zero inference-time cost?

We make that failure mode measurable, and then substantially mitigate it. The design principle is \textbf{attribution by freezing}: the CONCH image and text encoders \cite{lu2024conch} are frozen; the diagnosis head is a fixed cosine rule against class-text embeddings; there is no trained classifier. The only trainable object is a hierarchical evidence-selection policy. Consequently, when methods are compared on the same data and label space, differences in accuracy, leakage, and evidence choice are attributable to the selection policy alone; temporal class-incremental changes additionally reflect the expanding label space, which is why we also measure selection drift directly.

Prior LoRA-based CL---Online-LoRA \cite{wei2025onlinelora}, InfLoRA \cite{liang2024inflora}, SD-LoRA \cite{wu2025sdlora}, CL-LoRA \cite{he2025cllora}, FlexEM \cite{chen2026flexem}---preserves classifier or backbone representations. We instead confine continual learning to the evidence-selection policy over a frozen diagnostic pathway, which lets us measure directly whether a model forgets \emph{where to look for evidence}.

\textbf{Contributions.}
(1)~A formulation of continual evidence selection with complete attribution: frozen backbone, frozen head, selection-only training, and selection-level forgetting metrics (selection Jaccard, utility retention) alongside standard accuracy and forgetting.
(2)~A task-free consolidated selector: per-task rank-4 LoRA residuals folded exactly into one shared policy at task boundaries; inference uses a single model, no task ID, no memory, no module retrieval, at constant cost.
(3)~A mechanism separation supported by controls: merging is a substrate, not a demonstrated defense---LoRA-and-merge alone shows no systematic preservation benefit over sequential fine-tuning (3/5 paired seeds)---while preservation comes from replay, dual-level selection distillation, and a utility hinge operating on stored \emph{behavioral snapshots} rather than stored images.
(4)~An evaluation with pre-registered gates and win-count statistics over five seeds, including honest reporting of removed components, of comparisons that did not reach systematicity, and of one pre-registered expectation our data refuted.

\begin{figure*}[t]
  \centering
  \includegraphics[width=\textwidth]{figures/fig1_architecture.png}
  \caption{\textbf{Continual hierarchical evidence selection.}
  \emph{Left, middle (within a task):} patches from a frozen CONCH image encoder are assigned to $J{=}8$ tissue groups by argmax cosine similarity to fixed tissue prompts; a group selector $\Fg$ scores groups, a largest-remainder allocation converts scores into integer quotas $\sum_j b_j = B$ in one shot, and a patch selector $\Fp$ picks the top $b_j$ patches within each group. Selected evidence is pooled (score-weighted, $\ell_2$-normalized) and classified by a \emph{frozen} cosine rule against class-text embeddings---no trained classifier exists anywhere in the pipeline.
  \emph{Right (across tasks):} each task trains only rank-4 LoRA adapters on $\Fg,\Fp$ (fire); at the boundary the residual is folded exactly into the shared selector, which is the only model used at inference (no task ID, no memory, no retrieval). A 512-entry Selection Memory stores behavioral snapshots---group/patch scores, utilities, and replay indices; \emph{no patches or features}---and drives selection distillation, utility preservation, and replayed diagnostic training. Base MLPs are frozen within a task and updated \emph{only} by the boundary merge.}
  \label{fig:arch}
\end{figure*}

\section{Related work}
\label{sec:related}
\textbf{Evidence selection in WSIs.} Multiple-instance learning aggregates patch features with learned attention \cite{ilse2018abmil,lu2021clam,shao2021transmil}, implicitly ranking evidence but entangling selection with a trained classifier. We make selection explicit and budgeted, and remove the trained classifier entirely so that selection is the sole learned behavior.

\textbf{LoRA-based continual learning.} LoRA \cite{hu2022lora} has become the vehicle of choice for parameter-efficient CL. Online-LoRA \cite{wei2025onlinelora} expands adapters on detected distribution shift; FlexEM \cite{chen2026flexem} maintains a module pool with dual-stage retrieval at inference; InfLoRA \cite{liang2024inflora} and O-LoRA \cite{wang2023olora} constrain per-task subspaces; CL-LoRA \cite{he2025cllora} splits task-shared and task-specific adapters; SD-LoRA \cite{wu2025sdlora} decouples magnitude and direction and, like us, deploys the final model directly without per-task component selection. These methods differ in \emph{when} adaptation is consolidated. Pool-and-retrieve methods keep modules separate and select at inference; deferred-composition methods combine independently trained task vectors post hoc \cite{ilharco2023task,yadav2023ties,marczak2024magmax}; we consolidate \emph{incrementally}: each residual is trained on top of the already-merged selector under preservation constraints, then folded in exactly. Our study object also differs: prior work preserves classifier or backbone representations, whereas we preserve a selection \emph{policy} over a frozen pathway.

\textbf{Classical CL families.} We compare against representatives of rehearsal \cite{lopezpaz2017gem,buzzega2020derpp,rebuffi2017icarl} and functional regularization \cite{li2018lwf} adapted to the selection setting; parameter-importance methods \cite{kirkpatrick2017ewc} are complementary and left to future work (Sec.~\ref{sec:limitations}).

\section{Method}
\label{sec:method}
\subsection{Problem setup and the attribution principle}
Tasks $\tau = 1,\dots,T$ arrive sequentially; task $\tau$ contributes WSIs with labels from its own class set (two classes per organ task in our benchmark). For a slide with patch features $\{x_i\}_{i=1}^{N}\!\subset\!\mathbb{R}^{512}$ (frozen CONCH image encoder, $256{\times}256$ patches at $10\times$, pre-extracted), the system must select a set $\mathcal{P}$ with $|\mathcal{P}| = B = 8$ and predict from $\mathcal{P}$ alone. Prediction is a fixed rule: score-weighted pooling of $\mathcal{P}$, $\ell_2$ normalization, cosine similarity to frozen class-text embeddings $\{t_c\}$, softmax over the classes under evaluation. \textbf{Training is task-aware; inference is task-free}: at test time the model receives no task identity and is evaluated class-incrementally over the union of all classes seen so far, as well as task-incrementally within each task's class pair. Because all non-selector parameters are frozen, cross-method differences evaluated on the same data and label space are attributable to selector updates; temporal class-incremental changes additionally reflect the expanding evaluation label space. We therefore also report \emph{leakage}, the fraction of predictions falling in the wrong task's class set, and complement it with direct selection-drift metrics that bypass the head entirely.

\subsection{Hierarchical selection within a task}
\label{ssec:hier}
Patches are grouped by semantics, not clustering: each $x_i$ is assigned to $\arg\max_j \cos(x_i, p_j)$ over $J{=}8$ fixed tissue prompts (tumor, stroma, necrosis, \dots, normal) encoded once by the frozen text encoder; group prototypes are means, $g_j = \mathrm{mean}(x_i \in G_j)$. Two structurally identical two-layer MLPs, $\mathrm{Linear}(1537{\to}256)\to\mathrm{GELU}\to\mathrm{Linear}(256{\to}1)$, implement the policy: the group selector $\Fg$ maps $g_j$ to a score $r_j$; a largest-remainder (Hamilton) allocation converts $\mathrm{softmax}(r)$ into integer quotas $b_j$ with $\sum_j b_j = B$ in a single shot; the patch selector $\Fp$ scores patches, and the top $b_j$ within each group are selected. Each selector operates on its active 512-d visual input ($g_j$ for $\Fg$, $x_i$ for $\Fp$), with fixed zero padding retained only for implementation compatibility with the module's original 1537-d input shape; the discarded task-query and state-conditioning variants behind that shape are described in Sec.~\ref{sec:removed}. The two MLPs total $787{,}970$ parameters.

Within-task training minimizes $\mathcal{L}_{\mathrm{ev}} = \mathcal{L}_{\mathrm{diag}} + \beta_s\mathcal{L}_{\mathrm{sem}} + \beta_u\mathcal{L}_{\mathrm{util}}$: the frozen-head diagnostic loss on the selected evidence; a patch-wise semantic prior aligning patch scores with prompt affinities; and a differentiable utility term rewarding patches with positive counterfactual gain. Selection is made differentiable through the score-weighted pooling; $\beta_s=\beta_u=0.1$ throughout.

\subsection{Continual protocol: exact incremental consolidation}
\label{ssec:cl}
\textbf{Per-task residuals, one shared selector.} For every task, rank-4 LoRA adapters \cite{hu2022lora} are attached to each linear layer of $\Fg$ and $\Fp$ (four layers; $16{,}400$ trainable parameters per task, ${\approx}2.1\%$ of the selector; $\alpha{=}r$). During task $\tau$ the base is frozen and only the adapters train; the forward pass materializes $W_{\mathrm{eff}} = W + BA$ at every step, so folding the residual in at the boundary,
\begin{equation}
W_\tau \;=\; W_{\tau-1} + \Delta W_\tau,\qquad \Delta W_\tau = B_\tau A_\tau,
\end{equation}
is an \emph{identity rewrite}: the network computes the same function immediately before and after the merge. The adapter is then re-initialized for the next task. Inference uses the single merged selector: no task ID, no adapter bank, no retrieval, constant parameter count.

\textbf{Task-1 protocol, stated explicitly.} In the main configuration the very first task also trains only LoRA, on a frozen \emph{randomly initialized} base. Two observations contextualize this choice without settling it: the output layer's residual is a $1{\times}256$ matrix of rank at most one, hence unconstrained by the rank-4 budget; and the merge-only substrate (A2) reaches final task-IL parity with full-parameter sequential fine-tuning (A1) ($0.8105$ vs.\ $0.8086$; Table~\ref{tab:main})---though this compares whole four-task lifecycles and does not isolate the task-1 initialization. Whether warm-starting the base is preferable is an empirical question; we evaluate it directly with the pre-registered warm-start ablation (task~1 trains the full base, LoRA thereafter) under a two-sided decision rule; the ablation found no systematic class-IL difference (warm start better in 2/5 seeds; Sec.~\ref{sec:ablations}), so the parameter-efficient LoRA-from-task-1 protocol is retained.

\textbf{Merging is a substrate, not a mechanism.} LoRA merging alone does not prevent forgetting: with adapters and boundary merges but no preservation losses, final class-incremental accuracy shows no systematic difference from naive sequential fine-tuning ($0.4466$ vs.\ $0.4774$; 3/5 seeds, wide seed spread). What merging buys is a fixed-cost sequential \emph{substrate}---a single consolidated policy with task-free inference. Preservation comes from the training-time losses of Sec.~\ref{ssec:mem}. Two boundary-related equivalences preempt confusion. First, ``merge at each boundary'' and ``keep earlier adapters frozen in place and merge once at the end'' compute the same function and, with matched initialization, the same training trajectory; we verify this to numerical tolerance with a unit test and spend no experiments on it. Second, composing \emph{independently trained} task residuals post hoc (task arithmetic \cite{ilharco2023task}) is \emph{not} the same algorithm: each residual then never conditions on prior knowledge, and the composed model never receives a training signal. We evaluate this as a separate family in Sec.~\ref{sec:ablations}.

\subsection{Selection Memory and preservation losses}
\label{ssec:mem}
The memory $\mathcal{M}$ ($|\mathcal{M}|\le 512$, global reservoir sampling---uniform over all stored experiences) holds \emph{behavioral snapshots}: per slide, the task tag, slide ID, terminal selection state, group scores $r_{\mathrm{old}}$, candidate indices, patch scores $s_{\mathrm{old}}$, and counterfactual utilities $u_{\mathrm{old}}$. \emph{No patches and no features are stored}; features are reloaded by slide ID when a snapshot is replayed. Because entries are sampled uniformly over the stream, each past task's share is proportional to what it contributed; capacity is deliberately small, and Sec.~\ref{sec:results} shows the mechanism matters more than the capacity.

During task $\tau$, each optimizer step draws one stored slide ($\mathrm{replay}\_k{=}1$, i.e., a 1:1 old:new ratio) and adds
\begin{equation}
\mathcal{L}_{CL} = \lambda_{kd}\big[\mathrm{KL}(r_{\mathrm{old}}\Vert r_{\mathrm{new}}) + \mathrm{KL}(s_{\mathrm{old}}\Vert s_{\mathrm{new}})\big] + \lambda_{eq}\max(0, U_{\mathrm{old}} - U_{\mathrm{new}}) + \lambda_{r}\,\mathcal{L}_{\mathrm{diag}}^{\mathrm{rep}},
\end{equation}
with $\lambda_{kd}{=}\lambda_{eq}{=}\lambda_{r}{=}1$. Three design points deserve emphasis. The distillation teacher is \emph{not a second network}: targets are the stored score snapshots, so no teacher forward pass or checkpoint is kept. The hinge penalizes only utility \emph{regression}, leaving improvement unpenalized. And the replay term is the ordinary task loss evaluated on replayed slides---replay is a data mechanism, not a new objective; all three terms join the current sample's loss in a \emph{single} backward pass and optimizer step, so replay arms take no extra gradient steps. With replay, distillation, and utility preservation disabled, the continual pipeline reduces exactly to the merge-only baseline (A2). At inference the memory is not consulted at all.

\section{Experimental setup}
\label{sec:setup}
\textbf{Benchmark.} Four TCGA organ tasks---ESCA, RCC, BRCA, LUNG---each a binary subtype problem; the primary sequence is ESCA${\to}$RCC${\to}$BRCA${\to}$LUNG. Train/test: $2{,}273$ slides train; test $n=15/76/93/95$ per task. Class-incremental evaluation uses the union of all classes seen so far; task-incremental uses each task's own pair; leakage is the fraction of predictions in the wrong task's class set. All runs: 5 epochs/task, lr $10^{-3}$, $B{=}8$, rank 4, seeds 0--4.

\textbf{Arms.} Per-task specialists (R1; oracle reference), joint offline training (R2), sequential full fine-tuning (A1), LoRA with boundary merging only (A2), successively adding replay (A3) and dual-level distillation (A4), the full method (A5; A4 + utility hinge), and a single continual adapter that is never reset (L2). The LoRA arms (A2--A5) share the same parameterization, task order, and within-task objective, differing only in the enabled preservation components; sequential fine-tuning (A1) is a full-parameter baseline, and R1/R2 are specialist and offline references.

\textbf{Metrics.} Final class-IL accuracy ($A_{\mathrm{Final}}$), final task-IL accuracy, leakage; Forgetting $F = \frac{1}{T-1}\sum_{j<T}[\max_{s\ge j} a_{s,j} - a_{T,j}]$ and Plasticity $\frac{1}{T}\sum_j a_{j,j}$ from the full accuracy matrix; and two selection-level metrics computed on old-task slides: \emph{Selection Jaccard} between the evidence set chosen right after the task was learned and the set chosen at the end of the sequence, and \emph{utility retention} $\Delta U = U_{\mathrm{final}} - U_{\mathrm{own}}$ (values reported in the text).

\textbf{Statistical discipline.} With $n{=}5$ seeds we report mean$\pm$sd and \emph{win counts} over seed-paired comparisons; we call a difference \emph{systematic} only at 5/5 and \emph{directional} at 4/5, and we do not report $p$-values at this seed count. All gates below were registered before the corresponding runs.

\begin{table*}[t]
\centering\footnotesize
\caption{Continual-learning main table (mean$\pm$sd over 5 seeds; class-IL unless noted; flat selector substrate---the canonical main-table configuration, with hierarchical-policy contrasts reported in the text). Rows form a lifecycle progression: full fine-tuning $\to$ PEFT substrate $\to$ memory $\to$ behavioral preservation. Dashes: not applicable. Leakage: lower is better.}
\label{tab:main}
\begin{tabular}{lcccccc}
\toprule
Method & $A_{\mathrm{Final}}\!\uparrow$ & Forgetting$\downarrow$ & Plasticity$\uparrow$ & Sel.\ Jaccard$\uparrow$ & task-IL$\uparrow$ & Leakage$\downarrow$ \\
\midrule
Per-task specialist (R1) & $0.8777{\pm}.020$ & --- & $0.8777$ & --- & $0.9027{\pm}.019$ & --- \\
Joint (R2) & $0.7789{\pm}.028$ & --- & $0.7789$ & --- & $0.8498{\pm}.029$ & --- \\
\midrule
SeqFT (A1) & $0.4774{\pm}.103$ & $0.5193$ & $0.8669$ & $0.0023$ & $0.8086{\pm}.061$ & $0.4408$ \\
Seq-LoRA, one adapter (L2) & $0.7885{\pm}.015$ & $0.0982$ & $0.8555$ & $0.0661$ & $0.8941{\pm}.018$ & $0.1254$ \\
LoRA-Merge (A2) & $0.4466{\pm}.109$ & $0.5798$ & $0.8814$ & $0.0010$ & $0.8105{\pm}.036$ & $0.4756$ \\
\;+ Replay (A3) & $0.7778{\pm}.015$ & $0.1102$ & $0.8514$ & $0.0864$ & $0.9073{\pm}.018$ & --- \\
\;+ Distillation (A4) & $0.7972{\pm}.046$ & $0.1063$ & $0.8689$ & $0.1752$ & $0.8969{\pm}.019$ & --- \\
\textbf{PathSelect (A5)} & $\mathbf{0.8239{\pm}.028}$ & $\mathbf{0.0539}$ & $0.8515$ & $0.1613$ & $\mathbf{0.9147{\pm}.014}$ & $\mathbf{0.1005}$ \\
\bottomrule
\end{tabular}
\end{table*}

\section{Main results}
\label{sec:results}
\textbf{Accuracy.} The full method gains $+34.6$ class-IL points over SeqFT ($0.8239$ vs.\ $0.4774$; 5/5 seeds), $+10.6$ task-IL points (5/5), and cuts leakage by $34.0$ points (5/5), closing most of the gap to the per-task oracle and surpassing joint training's final accuracy. The merge-only substrate (A2) shows no systematic difference from sequential fine-tuning (A1) in the paired five-seed comparison (3/5)---the substrate result anticipated in Sec.~\ref{ssec:cl}.~\ref{ssec:cl}.

\textbf{Forgetting is halved, plasticity is untouched.} The full method's (A5) class-IL forgetting is $0.0539$ vs.\ $0.1102$ for replay alone (A3); its task-IL forgetting is $0.0080$. Plasticity remains within a narrow $0.85$--$0.88$ range across arms; the large retention gains are not accompanied by an obvious collapse in new-task learning.

\textbf{Forgetting where to look, measured directly.} Selection Jaccard separates the field by two orders of magnitude: unprotected arms re-select almost nothing of their original evidence ($0.0023$, $0.0010$), replay recovers some ($0.0864$), and distillation-based arms retain the most ($0.1752$, $0.1613$). Exact evidence-set overlap and diagnostic performance need not move monotonically together: A4 retains slightly higher Jaccard, whereas A5 achieves better final accuracy and forgetting; we report both. The revised utility-retention metric orders arms consistently with accuracy: $\Delta U$ is $-16.8$ for the full method and $-21.2$/$-31.9$ for replay/distillation, versus $-220$/$-301$ for the unprotected sequential baselines (A1/A2) (single adapter L2: $-27.5$; warm start W1 best overall at $-12.8$)---every arm loses old-task utility, but protected arms lose an order of magnitude less. We also note honestly that the A5$-$A4 margin flips sign under some task orders; the class-IL gain of the utility hinge is order-sensitive and we present it as such.

\textbf{Hierarchy and dual-level distillation.} Under the flat substrate the full-method-over-replay margin (A5$-$A3) on task-IL is $+0.74{\pm}1.93$ (3/5); the hierarchical policy strengthens the same contrast to $+3.28{\pm}2.40$ task-IL and $+5.76{\pm}3.42$ class-IL (both 5/5), while its \emph{absolute} A5 accuracy stays within noise of flat ($-0.58$/$-1.20$, 1/5) with somewhat higher leakage ($+2.15$, 0/5)---the architecture was adopted under pre-registered criteria for amplifying and stabilizing the method contrast, and we report both sides. Removing group-level distillation under the hierarchy costs $3.95$ task-IL and $3.71$ class-IL points (both 5/5), supporting an additional contribution from group-level preservation on top of the patch-level pathway.

\textbf{Preservation design can outweigh buffer size.} In this controlled comparison on the hierarchical configuration, the full method with a 128-entry memory outperforms replay-only with 512 entries ($+3.06$ task-IL, 5/5), indicating that \emph{what} is stored---behavioral snapshots consumed by distillation and the hinge---can matter more than how much.

\section{Update-protocol ablations (pre-registered)}
\label{sec:ablations}
All arms below were registered with two-sided decision rules before any run; win counts over 5 seeds decide, and the incumbent is retained on a split. One pre-registered expectation was refuted; we report it as measured.

\textbf{Warm start (W1).} Task 1 trains the full base; LoRA thereafter. Class-IL shows no systematic difference (W1 $-1.80$pp mean, better in 2/5 seeds; task-IL $-0.96$, 1/5), so the parameter-efficient incumbent is retained under the pre-registered split rule. Warm-starting does reshape selection: Jaccard rises to $0.245$ ($+0.06$ over the full method, 4/5) and $\Delta U$ is the best of any arm ($-12.8$)---changing \emph{what} is selected without changing outcomes, echoing the removed-components pattern. Base quality matters precisely when nothing else protects: without preservation losses, its bare variant (W1B) improves over the bare substrate (A2) by $+2.86$ class-IL (4/5) with $-1.59$ leakage (4/5); with them, the advantage dissolves---preservation subsumes initialization.

\textbf{One adapter vs.\ fresh-per-task (L2 vs.\ A5).} Skipping the boundary reset caps the total drift at rank 4; fresh adapters allow rank up to 16. On the flat substrate the capacity matters: the single-adapter variant loses $3.54$ class-IL (full method better in 4/5) and $2.06$ task-IL (4/5) points, supporting the fresh-adapter design the merge protocol assumes. (Bare, the single adapter (L2B) edges the substrate (A2) by $+1.37$ task-IL, 4/5---mild continuity benefits exist, but capacity dominates once preservation is active.) On the hierarchical substrate the same paired comparison becomes a trade-off rather than a verdict: fresh adapters remain systematically better at within-task discrimination ($+1.25$ task-IL, 5/5), while the single adapter leaks systematically less ($-3.25$, 5/5) and is directionally ahead on class-IL (4/5)---the hierarchy's quota structure may absorb part of the load that adapter capacity carries on the flat substrate. The required capacity thus interacts with the selection architecture; we retain the fresh-adapter design (incumbent, and favored on a primary axis) and leave a registered study of this interaction to the extended version.

\textbf{Post-hoc composition (C1 sum, C2 average).} The pre-registered expectation was refuted---informatively. Residuals trained independently from a shared $\theta_0$ and then summed beat the bare sequential substrate (A2) on every metric in every seed ($+24.6$ class-IL, $+6.6$ task-IL, $-24.9$ leakage; all 5/5): without preservation, sequential updating is actively destructive---each task's update freely overwrites its predecessors---whereas independent residuals never interfere during training and compose well, echoing the task-arithmetic literature \cite{ilharco2023task}. Composition is not the endpoint: its defining independence excludes replay, distillation targets require an accumulated model, and the full method beats it decisively ($+13.1$ class-IL and $-12.6$ leakage, both 5/5; $+3.8$ task-IL, 4/5). Sum and average are interchangeable here (no systematic difference). The refined claim: \emph{sequential consolidation earns its place only in combination with behavioral preservation---the component that requires sequence to define}. Composition also attains the highest selection Jaccard of any arm ($0.275$) while trailing the full method by 13 accuracy points---a second instance, after warm start, of interventions that stabilize \emph{what is selected} without improving \emph{what is diagnosed}.

\textbf{Damped merging (A5H, $\alpha{=}0.5$).} Halving each fold-in leaves class-IL within noise (2/5); full merging is directionally better on task-IL ($+0.98$, 4/5) and leakage ($-1.15$, 4/5), so $\alpha{=}1$ is retained. Damping raises Jaccard ($0.238$ vs.\ $0.161$)---the third such instance: selection stabilizes, accuracy does not.

\textbf{Replay budget.} Sensitivity to the replay draw and to memory-policy variants is deferred to the extended version.

\section{Removed components}
\label{sec:removed}
Three architectural components failed their pre-registered gates and were removed from the method; we report them because removal is not hiding. In the general parameterization the selectors accept $[\,\mathrm{feat};q_\tau;e;\tilde b\,]\in\mathbb{R}^{1537}$; the main configuration zero-fills the conditioning blocks after the corresponding variants failed their gates. The task-conditioning query $q_\tau$ failed its gate; separately, task identity is already highly decodable from slide-average features ($98.21\%$) and group prototypes ($98.57\%$), consistent with---though not proving---the hypothesis that explicit conditioning adds little new information on this benchmark. It nonetheless carries a reproducible secondary effect, reducing leakage by $5.92$ points (5/5), which we retain as an observation, not a mechanism. Group-level semantic priors and stateful (sequential) selection likewise failed their gates; the group-level prior carries its own secondary finding, reducing quota-distribution divergence by $0.005$ ($-0.005{\pm}0.004$, 5/5) without translating into accuracy. Full gate cards appear in the extended version.

\section{Limitations and discussion}
\label{sec:limitations}
Five seeds and win-count discipline bound what we can claim; we report no $p$-values by design. Task boundaries are known during training (organ cohorts arrive staged), which is clinically motivated but weaker than task-free online CL; our inference, however, is fully task-free. We compare within parameter-isolation, rehearsal, and functional-regularization families; parameter-importance methods (e.g., EWC) are complementary and deferred. Results use one foundation model (CONCH) and one budget ($B{=}8$); the removed-component analysis suggests our benchmark's tasks are unusually separable, and harder mixtures may re-open the case for explicit task conditioning.

\section{Conclusion}
\label{sec:conclusion}
Freezing everything except a budgeted selection policy turns continual learning for WSIs into a question that can be asked cleanly and answered mechanistically: models do forget where to look, merging alone does not stop them, and training-time behavioral preservation---replay, dual-level selection distillation, and a utility hinge over stored snapshots---substantially mitigates it, at zero inference-time cost.

% To start a new column (but not a new page) and help balance the last-page
% column length use \vfill\pagebreak.
% -------------------------------------------------------------------------
%\vfill
%\pagebreak

\bibliographystyle{IEEEbib}
\bibliography{refs}

\end{document}
